
%%%%%%%%%%%%%%%%%%%%%%%%%%%%%%%%%%%%%%%%%%%%%%%%%%%%%%%%%%%%%%%%%%%%%%%%%%%%%%%%%%
\begin{frame}[fragile]\frametitle{}
\begin{center}
{\Large Basic Prompt Engineering Techniques}
\end{center}
\end{frame}



%%%%%%%%%%%%%%%%%%%%%%%%%%%%%%%%%%%%%%%%%%%%%%%%%%%%%%%%%%%
\begin{frame}[fragile]\frametitle{Basic Prompt vs. Engineered Prompt}

\begin{columns}
    \begin{column}[T]{0.5\linewidth}
		\begin{center}
		\includegraphics[width=\linewidth,keepaspectratio]{promptengg64}
		{\tiny Basic prompt: generates a list of product name ideas.}
		\end{center}	
    \end{column}
    \begin{column}[T]{0.5\linewidth}
		\begin{center}
		\includegraphics[width=\linewidth,keepaspectratio]{promptengg65}
		{\tiny Engineered prompt: adds role, constraints, and format.}
		\end{center}
    \end{column}
  \end{columns}

\medskip
\textbf{Lesson:} The engineered prompt is more specific: it defines the \emph{target audience}, the \emph{style}, and the \emph{output format}. Same AI, much better result.

{\tiny (Ref: The Complete Prompt Engineering for AI Bootcamp (2023))}

\end{frame}


%%%%%%%%%%%%%%%%%%%%%%%%%%%%%%%%%%%%%%%%%%%%%%%%%%%%%%%%%%%
\begin{frame}[fragile]\frametitle{Elements of a Prompt}

A well-crafted prompt can contain any combination of:

\begin{center}
\includegraphics[width=0.5\linewidth,keepaspectratio]{promptengg28}

{\tiny (Ref: Prompt Engineering Overview -- Elvis Saravia)}
\end{center}		

\textbf{In plain words:}
\begin{itemize}
\item \textbf{Instruction}: What do you want done? (\emph{``Summarise this text''})
\item \textbf{Context}: Background information to guide the answer
\item \textbf{Input Data}: The actual content to process
\item \textbf{Output Indicator}: The format you want (\emph{``in bullet points'', ``as JSON''})
\end{itemize}
		
\end{frame}

%%%%%%%%%%%%%%%%%%%%%%%%%%%%%%%%%%%%%%%%%%%%%%%%%%%%%%%%%%%
\begin{frame}[fragile]\frametitle{Length, Tone, Style, Audience, and Context Control}

You can control almost every aspect of the AI's output just by adding a few words to your prompt:

\medskip
\begin{center}
\renewcommand{\arraystretch}{1.4}
\begin{tabular}{|l|p{5.5cm}|}
\hline
\textbf{Control Type} & \textbf{Example Prompt Addition} \\
\hline
Length & \emph{``in exactly 100 words''} \\
Tone & \emph{``in a friendly, encouraging tone''} \\
Style & \emph{``written like a newspaper editorial''} \\
Audience & \emph{``explain this to a 10-year-old''} \\
Context & \emph{``given that the user is a first-time investor''} \\
\hline
\end{tabular}
\end{center}

\medskip
\textbf{Pro tip:} Combine multiple controls in one prompt. They do not conflict: the AI respects all of them simultaneously.

{\tiny (Ref: Prompt Engineering, Sudalai Rajkumar)}

\end{frame}

%%%%%%%%%%%%%%%%%%%%%%%%%%%%%%%%%%%%%%%%%%%%%%%%%%%%%%%%%%%
\begin{frame}[fragile]\frametitle{The Five Pillars of Prompting}

\begin{itemize}
\item \textbf{Providing Examples}: Show the AI what a good answer looks like (Few-Shot)
\item \textbf{Giving Directions}: Tell the AI exactly what role, task, and constraints to follow
\item \textbf{Formatting Responses}: Specify the output structure (JSON, table, bullet list, etc.)
\item \textbf{Evaluating Quality}: Test prompts systematically; iterate based on failures
\item \textbf{Chaining AIs}: Feed the output of one AI call as input to another for complex tasks
\end{itemize}	

\medskip
All advanced techniques are variations or combinations of these five fundamentals.

\end{frame}

%%%%%%%%%%%%%%%%%%%%%%%%%%%%%%%%%%%%%%%%%%%%%%%%%%%%%%%%%%%
\begin{frame}[fragile]\frametitle{Providing Examples in Prompts}

\begin{columns}
    \begin{column}[T]{0.55\linewidth}
		\begin{center}
		\includegraphics[width=\linewidth,keepaspectratio]{promptengg66}

		{\tiny (Ref: The Complete Prompt Engineering for AI Bootcamp (2023))}
		\end{center}	
    \end{column}
    \begin{column}[T]{0.45\linewidth}
    	\textbf{Why show examples?}
    	
    	\medskip
		Showing the AI what a good answer looks like is often \emph{more effective} than describing it in words.

		\medskip
		Give \textbf{diverse} examples to avoid the AI copying only one style.

		\medskip
		Think of it like showing a student sample exam papers before the test: they immediately understand the format and expectation.
    \end{column}
  \end{columns}
\end{frame}


%%%%%%%%%%%%%%%%%%%%%%%%%%%%%%%%%%%%%%%%%%%%%%%%%%%%%%%%%%%
\begin{frame}[fragile]\frametitle{Prompting by Examples (Few-Shot in Action)}

\begin{center}
\includegraphics[width=\linewidth,keepaspectratio]{promptengg8}
{\tiny (Ref: Cohere https://txt.cohere.ai/generative-ai-part-1/)}
\end{center}		
\end{frame}


%%%%%%%%%%%%%%%%%%%%%%%%%%%%%%%%%%%%%%%%%%%%%%%%%%%%%%%%%%%
\begin{frame}[fragile]\frametitle{Zero-Shot Prompting}
    \begin{itemize}
        \item \textbf{What it is:} You ask the AI to perform a task without giving it \emph{any} examples first. You just give the instruction and the input.
        \item \textbf{When it works:} Simple, well-defined tasks that the AI has seen many examples of during training.
        \item \textbf{When it struggles:} Complex reasoning, niche topics, or tasks requiring a very specific output format.
		\item \textbf{Intuition:} You are trusting the AI's general knowledge without coaching it. Like asking a stranger a question: they may get it right, but you cannot guarantee it.
    \end{itemize}
	

	
\begin{lstlisting}[language=HTML]
Prompt:
Classify the text into neutral, negative, or positive. 

Text: I think the vacation is okay.
Sentiment:

Output:
Neutral
\end{lstlisting}

\end{frame}


%%%%%%%%%%%%%%%%%%%%%%%%%%%%%%%%%%%%%%%%%%%%%%%%%%%%%%%%%%%
\begin{frame}[fragile]\frametitle{Few-Shot Prompting}
    \begin{itemize}
        \item \textbf{What it is:} You give the AI a few examples (``shots'') of input $\rightarrow$ desired output before posing your actual question. The AI learns the \emph{pattern} from your examples.
        \item \textbf{When it works:} Whenever you need consistent output format or style.
		\item \textbf{Intuition:} You are teaching by example rather than by description. The AI picks up the pattern and continues it.
    \end{itemize}



        \begin{lstlisting}[language=HTML]
A "whatpu" is a small, furry animal native to Tanzania. 

An example sentence using "whatpu":
We saw these very cute whatpus in Africa.

To "farduddle" means to jump up and down really fast. 

An example sentence using "farduddle":

Model output:
When we won the game, we all started to farduddle!
        \end{lstlisting}


\end{frame}


%%%%%%%%%%%%%%%%%%%%%%%%%%%%%%%%%%%%%%%%%%%%%%%%%%%%%%%%%%%
\begin{frame}[fragile]\frametitle{Limitations of Few-Shot Prompting}
    \begin{itemize}
        \item Few-shot prompting excels at format and style: but it \textbf{does not help with reasoning}.
        \item If the AI needs to \emph{think through} a problem, simply seeing examples is not enough.
		\item \textbf{Fix:} Use Chain-of-Thought prompting (next section).
    \end{itemize}
	
        \begin{lstlisting}[language=HTML]
The odd numbers in this group add up to an even number: 
4, 8, 9, 15, 12, 2, 1. ,> A: False.
17, 10, 19, 4, 8, 12, 24,> A: True.
16, 11, 14, 4, 8, 13, 24,> A: True.
17,  9, 10, 12, 13, 4, 2,> A: False.

The odd numbers in this group add up to an even number: 
15, 32, 5, 13, 82, 7, 1. 
A:

Model output:
The answer is True.  [WRONG! Correct answer is False]
\end{lstlisting}



\end{frame}


%%%%%%%%%%%%%%%%%%%%%%%%%%%%%%%%%%%%%%%%%%%%%%%%%%%%%%%%%%%
\begin{frame}[fragile]\frametitle{Giving Direction }

\begin{columns}
    \begin{column}[T]{0.5\linewidth}
		\begin{center}
		\includegraphics[width=\linewidth,keepaspectratio]{promptengg69}
		{\tiny Direction: describe what you are imagining.}
		\end{center}

		{\tiny (Ref: The Complete Prompt Engineering for AI Bootcamp (2023))}
    \end{column}
    \begin{column}[T]{0.5\linewidth}
    \textbf{Direction tips:}
    \begin{itemize}
    \item Use keywords relevant to the desired output.
    \item Describe the scenario the answer should fit.
    \item ``Keyword-based'' and ``scenario-based'' guiding both work.
    \end{itemize}

    \end{column}
  \end{columns}
\end{frame}

%%%%%%%%%%%%%%%%%%%%%%%%%%%%%%%%%%%%%%%%%%%%%%%%%%%%%%%%%%%
\begin{frame}[fragile]\frametitle{Formatting Responses}

\begin{columns}
    \begin{column}[T]{0.5\linewidth}

		\begin{center}
		\includegraphics[width=\linewidth,keepaspectratio]{promptengg71}
		{\tiny Formatting: define the response structure.}
		\end{center}
		{\tiny (Ref: The Complete Prompt Engineering for AI Bootcamp (2023))}
    \end{column}
    \begin{column}[T]{0.5\linewidth}

    \textbf{Formatting tips:}
    \begin{itemize}
    \item Specify JSON, Markdown, table, or numbered list.
    \item For API integrations, JSON output makes the response directly parseable by code.
    \end{itemize}
    \end{column}
  \end{columns}
\end{frame}


%%%%%%%%%%%%%%%%%%%%%%%%%%%%%%%%%%%%%%%%%%%%%%%%%%%%%%%%%%%
\begin{frame}[fragile]\frametitle{Evaluating Quality}

\begin{columns}
    \begin{column}[T]{0.5\linewidth}
		\begin{center}
		
\includegraphics[width=\linewidth,keepaspectratio]{promptengg73}
		{\tiny Evaluating: test prompts systematically.}
		\end{center}	
	
		{\tiny (Ref: The Complete Prompt Engineering for AI Bootcamp (2023))}
    \end{column}
    \begin{column}[T]{0.5\linewidth}
    \textbf{Evaluating Quality:}
    \begin{itemize}
    \item Try different prompt variations systematically.
    \item Build a small set of test cases with known correct answers.
    \item Identify where the prompt fails and refine it.
    \end{itemize}
    

    \end{column}
  \end{columns}
\end{frame}

%%%%%%%%%%%%%%%%%%%%%%%%%%%%%%%%%%%%%%%%%%%%%%%%%%%%%%%%%%%
\begin{frame}[fragile]\frametitle{Chaining AIs}

\begin{columns}
    \begin{column}[T]{0.5\linewidth}

		\begin{center}
		\includegraphics[width=\linewidth,keepaspectratio]{promptengg75}
		{\tiny Chaining: output of one AI feeds into another.}
		\end{center}	
		{\tiny (Ref: The Complete Prompt Engineering for AI Bootcamp (2023))}
    \end{column}
    \begin{column}[T]{0.5\linewidth}
    \textbf{Chaining AIs:}
    \begin{itemize}
    \item Use the output of one AI call as the input for another.
    \item Example: AI 1 extracts key data $\rightarrow$ AI 2 formats it as a report $\rightarrow$ AI 3 translates it.
    \item Enables complex, multi-stage workflows with simple prompts.
    \end{itemize}
    \end{column}
  \end{columns}
\end{frame}


